Electricity consumption is a low-salience decision. Households receive the benefits of electricity use immediately, while the monetary and environmental costs are typically delayed, aggregated, and difficult to attribute to specific activities. Following \citet{alcocerSalienceBiasFramework2024}, we use a simple salience framework to motivate why digital feedback may reduce electricity consumption and to derive the empirical predictions tested below.

Consider a household that chooses electricity consumption \(e\) and a composite good \(x\). Under full salience, the household solves
\[
\max_{e,x} u(e,x)
\quad \text{s.t.} \quad
p_e e + p_x x \leq w,
\]
where \(p_e\) is the price of electricity, \(p_x\) is the price of the composite good, and \(w\) is income. Let the full-salience solution be \((e^*,x^*)\). If electricity costs are not fully salient, the household behaves as if the budgetary cost of electricity were lower than it actually is. Let \(\tau \in (0,1]\) denote the salience of electricity costs. The perceived budget constraint is then
\[
\widetilde F(x,e)=p_xx+\tau p_e e-w=0.
\]
When \(\tau<1\), the perceived marginal cost of electricity is below the true marginal cost. Under standard assumptions, this leads to higher electricity consumption than under full salience.

Figure~\ref{fig:Theory} illustrates this mechanism. Panel~(a) shows the full-salience benchmark. The household chooses the bundle \((e^*,x^*)\), where the indifference curve is tangent to the real budget constraint. Panel~(b) shows the low-salience case. If electricity costs are less salient, the household behaves as if electricity were cheaper than it actually is. The perceived, or behavioral, budget constraint therefore induces a higher level of electricity consumption. Once the real budget constraint binds, this higher electricity consumption leaves fewer resources for other goods. The figure captures the central intuition of the model: low salience can generate overconsumption of electricity relative to the full-information benchmark.

\begin{figure}[h]
    \centering
    \begin{subfigure}[b]{0.72\textwidth}
        \centering
        \includegraphics[width=\textwidth]{paper/figures/Theory_plot_a (1).png}
        \caption{Full salience}
        \label{fig:Theory_a}
    \end{subfigure}

    \vspace{0.4cm}

    \begin{subfigure}[b]{0.72\textwidth}
        \centering
        \includegraphics[width=\textwidth]{paper/figures/Theory plot b (1).png}
        \caption{Low salience}
        \label{fig:Theory_b}
    \end{subfigure}

    \caption{Salience bias and electricity consumption. Panel~(a) shows the full-salience benchmark, in which the household chooses \((e^*,x^*)\) on the real budget constraint. Panel~(b) shows the low-salience case, in which the household perceives electricity as less costly, chooses higher electricity consumption, and ends up with fewer resources for other goods once the real budget constraint binds. Adapted from \citet{alcocerSalienceBiasFramework2024}.}
    \label{fig:Theory}
\end{figure}

The smartphone application is designed to increase the salience of electricity use by making consumption, monetary costs, and emissions more visible. Feedback on recent electricity use reduces the delay between behavior and information, while cost and emissions feedback make the consequences of consumption easier to interpret. Social comparison and gamified learning features may further increase attention to electricity use. The framework therefore motivates our first research question: does access to the application reduce household electricity consumption relative to the control group? Under the salience interpretation, we expect households offered access to the application to consume less electricity than households without access.

The framework also guides two additional analyses. If the response reflects behavioral adjustment, reductions should be more likely where households have scope to change consumption. 

We therefore distinguish between daytime consumption (07:00–22:00), when discretionary adjustments are most feasible, and night-time consumption outside this window.

This distinction is an empirical proxy rather than an appliance-level classification, but it allows us to examine whether reductions are concentrated in more discretionary electricity use. This leads to the second research question: are treatment effects larger for flexible than for non-flexible consumption?

Finally, the model highlights the role of attention over time. If the application mainly works by making electricity costs temporarily more salient, treatment effects should attenuate as app usage declines. If households learn about their consumption patterns or form new habits, some reductions may persist even after engagement falls. The third research question therefore asks whether the time profile of treatment effects coincides with the time profile of app engagement. Because engagement is measured after treatment assignment, this analysis is descriptive and is not interpreted as the causal effect of app usage.